% ============================================================
%GUÍA DE FÍSICA IV — CAPÍTULOS 34 Y 35
% El Campo Magnético · La Ley de Ampère
% CECyT 9 IPN · Autor: CURIEL GARCIA JUAN JESUS
% ============================================================

\documentclass[12pt, letterpaper]{report}

% ── Codificación e idioma ────────────────────────────────────

\usepackage[spanish, mexico]{babel}

% ── Márgenes ────────────────────────────────────────────────
\usepackage[top=2.5cm, bottom=2.8cm, left=2.8cm, right=2.5cm]{geometry}

% ── Fuentes ─────────────────────────────────────────────────
\usepackage{lmodern}
\usepackage{microtype}

% ── Matemáticas ─────────────────────────────────────────────
\usepackage{amsmath, amssymb, amsthm, bm}
\usepackage{mathtools}

% ── Color y cajas ───────────────────────────────────────────
\usepackage[dvipsnames, table]{xcolor}
\usepackage{tcolorbox}
\tcbuselibrary{skins, breakable, theorems}

% ── Cabeceras / pies ────────────────────────────────────────
\usepackage{fancyhdr}
\usepackage{titlesec}
\usepackage{titletoc}

% ── Tablas y figuras ────────────────────────────────────────
\usepackage{booktabs}
\usepackage{array}
\usepackage{longtable}
\usepackage{multirow}
\usepackage{graphicx}
\usepackage{float}
\usepackage{caption}
\usepackage{mdframed}

% ── Hipervínculos ───────────────────────────────────────────
\usepackage[colorlinks=true, linkcolor=NavyBlue,
            urlcolor=Guinda, citecolor=NavyBlue,
            pdftitle={Guia Fisica IV - Caps 34 y 35},
            pdfauthor={Curiel Garcia Juan Jesus}]{hyperref}

% ── Numeración ──────────────────────────────────────────────
\numberwithin{equation}{section}
\usepackage{enumitem}
\usepackage{setspace}
\setstretch{1.15}

% ── Espaciado entre párrafos ─────────────────────────────────
\setlength{\parskip}{4pt}
\setlength{\parindent}{0pt}

% ============================================================
%  PALETA DE COLORES — AZUL MARINO Y GUINDA (IPN)
% ============================================================
\definecolor{NavyBlue}{RGB}{0, 32, 96}
\definecolor{Guinda}{RGB}{130, 0, 30}
\definecolor{LightNavy}{RGB}{210, 220, 240}
\definecolor{LightGuinda}{RGB}{245, 215, 220}
\definecolor{BoxGray}{RGB}{248, 248, 252}
\definecolor{GoldAccent}{RGB}{180, 140, 0}
\definecolor{DarkGray}{RGB}{50, 50, 60}

% ============================================================
%  FORMATO DE TÍTULOS
% ============================================================

\titleformat{\chapter}[display]
  {\normalfont\LARGE\bfseries\color{NavyBlue}}
  {\filleft
   \colorbox{Guinda}{\color{white}\quad
   \fontsize{48}{48}\selectfont\thechapter\quad}
   \vspace{2pt}\color{Guinda}\titlerule[3pt]}
  {0.5em}
  {\filright\color{NavyBlue}}
  [\vspace{4pt}\color{Guinda}\titlerule[1.5pt]\vspace{6pt}]

\titlespacing{\chapter}{0pt}{-20pt}{20pt}

\titleformat{\section}
  {\normalfont\large\bfseries\color{NavyBlue}}
  {\colorbox{NavyBlue}{\color{white}\;\thesection\;}}{0.7em}{}
  [\color{Guinda}\titlerule[0.8pt]\vspace{2pt}]

\titlespacing{\section}{0pt}{14pt}{8pt}

\titleformat{\subsection}
  {\normalfont\normalsize\bfseries\color{Guinda}}
  {\thesubsection\;}{0.5em}{}
  [\color{NavyBlue!30}\titlerule[0.4pt]]

\titlespacing{\subsection}{0pt}{10pt}{5pt}

\titleformat{\subsubsection}
  {\normalfont\normalsize\bfseries\color{NavyBlue!80}}
  {\thesubsubsection\;}{0.5em}{}

% ============================================================
%  CABECERAS Y PIES DE PÁGINA
% ============================================================
\pagestyle{fancy}
\fancyhf{}

\fancyhead[L]{%
  \small\color{NavyBlue}\textbf{\leftmark}}
\fancyhead[R]{%
  \small\color{Guinda}\textbf{Física IV · CECyT 9 IPN}}

\fancyfoot[L]{%
  \footnotesize\color{DarkGray}\textit{Curiel García Juan Jesús}}
\fancyfoot[C]{%
  \footnotesize\color{NavyBlue!60}\rule{0pt}{8pt}}
\fancyfoot[R]{%
  \small\color{NavyBlue}\textbf{\thepage}}

\renewcommand{\headrulewidth}{1.2pt}
\renewcommand{\headrule}{%
  \hbox to\headwidth{\color{NavyBlue}\leaders\hrule height\headrulewidth\hfill}}
\renewcommand{\footrulewidth}{0.6pt}
\renewcommand{\footrule}{%
  \hbox to\headwidth{\color{Guinda!50}\leaders\hrule height\footrulewidth\hfill}}

\fancypagestyle{plain}{%
  \fancyhf{}
  \fancyfoot[C]{\small\color{NavyBlue}\textbf{\thepage}}
  \renewcommand{\headrulewidth}{0pt}
  \renewcommand{\footrulewidth}{0.6pt}
}

% ============================================================
%  CAJAS TEMÁTICAS
% ============================================================
\tcbset{
  %% Caja de definición — azul marino
  defbox/.style={
    enhanced, breakable,
    colback=LightNavy!60, colframe=NavyBlue,
    fonttitle=\bfseries\small\color{white},
    coltitle=white, colbacktitle=NavyBlue,
    attach boxed title to top left={yshift=-2.5mm, xshift=5mm},
    boxed title style={colframe=NavyBlue, colback=NavyBlue,
      sharp corners, rounded corners=northeast},
    arc=3pt, sharp corners=south,
    top=5mm, left=4mm, right=4mm, bottom=4mm,
    boxrule=0.8pt
  },
  %% Caja de propiedad / teorema — guinda
  propbox/.style={
    enhanced, breakable,
    colback=LightGuinda!70, colframe=Guinda,
    fonttitle=\bfseries\small\color{white},
    coltitle=white, colbacktitle=Guinda,
    attach boxed title to top left={yshift=-2.5mm, xshift=5mm},
    boxed title style={colframe=Guinda, colback=Guinda,
      sharp corners, rounded corners=northeast},
    arc=3pt, sharp corners=south,
    top=5mm, left=4mm, right=4mm, bottom=4mm,
    boxrule=0.8pt
  },
  %% Caja de ejemplo / ejercicio — gris claro
  exbox/.style={
    enhanced, breakable,
    colback=BoxGray, colframe=NavyBlue!45,
    fonttitle=\bfseries\small\color{NavyBlue},
    coltitle=NavyBlue, colbacktitle=NavyBlue!12,
    attach boxed title to top left={yshift=-2.5mm, xshift=5mm},
    boxed title style={colframe=NavyBlue!45, colback=NavyBlue!12,
      sharp corners},
    arc=3pt,
    top=5mm, left=4mm, right=4mm, bottom=4mm,
    boxrule=0.6pt
  },
  %% Caja de nota / advertencia — amarillo tenue
  notebox/.style={
    enhanced, breakable,
    colback=yellow!7, colframe=GoldAccent,
    fonttitle=\bfseries\small\color{GoldAccent},
    coltitle=GoldAccent, colbacktitle=yellow!20,
    attach boxed title to top left={yshift=-2.5mm, xshift=5mm},
    boxed title style={colframe=GoldAccent, colback=yellow!20,
      sharp corners},
    arc=3pt,
    top=5mm, left=4mm, right=4mm, bottom=4mm,
    boxrule=0.6pt
  },
  %% Marco para figuras
  figbox/.style={
    enhanced,
    colback=BoxGray, colframe=NavyBlue!30,
    arc=4pt,
    top=8mm, bottom=8mm, left=6mm, right=6mm,
    boxrule=0.5pt
  }
}

% ============================================================
%  COMANDOS ÚTILES
% ============================================================
\newcommand{\vB}[1]{\boldsymbol{#1}}          % vector en negritas
\newcommand{\un}[1]{\,\mathrm{#1}}            % unidades
\newcommand{\ec}[1]{\textbf{(Ec.\;#1)}}       % referencia ecuación

%% Figura placeholder
\newcommand{\figplaceholder}[2]{%
  \begin{tcolorbox}[figbox]
    \centering
    \textcolor{NavyBlue!50}{\rule{0pt}{2.8cm}}\\[4pt]
    {\small\color{NavyBlue}\textbf{[FIGURA\ #1]}\\[2pt]
    \textit{#2}}
  \end{tcolorbox}
}

%% Ecuación destacada con recuadro de color
\newcommand{\ecimportante}[1]{%
  \begin{equation}
    \tcboxmath[colback=LightNavy!80, colframe=NavyBlue,
      boxrule=0.8pt, arc=2pt]{#1}
  \end{equation}
}

% ============================================================
\begin{document}
% ============================================================

% ──────────────────────────────────────────────────────────────
%  PORTADA
% ──────────────────────────────────────────────────────────────
\begin{titlepage}
  \pagecolor{NavyBlue}
  \color{white}
  \centering
  \vspace*{1.8cm}

  % Línea decorativa superior
  {\color{Guinda}\rule{0.55\textwidth}{2.5pt}}\\[0.6cm]

  {\fontsize{13}{16}\selectfont\scshape Instituto Politécnico Nacional}\\[0.25cm]
  {\large Centro de Estudios Científicos y Tecnológicos No.\,9}\\[0.1cm]
  {\normalsize\itshape ``Juan de Dios Bátiz Paredes''}\\[1.2cm]

  {\color{Guinda}\rule{0.55\textwidth}{1.5pt}}\\[1cm]

  {\fontsize{16}{20}\bfseries\selectfont GUÍA DE ESTUDIO}\\[0.4cm]
  {\fontsize{28}{34}\bfseries\selectfont Física IV}\\[0.6cm]

  {\color{Guinda}\rule{0.4\textwidth}{1pt}}\\[0.8cm]

  {\fontsize{18}{22}\bfseries\selectfont
   Capítulo 34 — El Campo Magnético}\\[0.4cm]
  {\fontsize{18}{22}\bfseries\selectfont
   Capítulo 35 — La Ley de Ampère}\\[1.5cm]

  {\color{Guinda}\rule{0.55\textwidth}{1.5pt}}\\[1.2cm]

  {\large Elaborado por:}\\[0.4cm]
  {\Large\bfseries CURIEL GARCÍA JUAN JESÚS}\\[2cm]

  \vfill

  {\small\color{white!65}
   Contenido basado en: Resnick, Halliday \& Krane —
   \textit{Física}, Vol.\,2 ·
   Documento compilado con \LaTeX}\\[0.2cm]
  {\small\color{white!65}\today}
\end{titlepage}

\pagecolor{white}
\color{black}

% ──────────────────────────────────────────────────────────────
%  TABLA DE CONTENIDOS
% ──────────────────────────────────────────────────────────────
\tableofcontents
\thispagestyle{fancy}
\clearpage

% ==============================================================
% ==============================================================
%  CAPÍTULO 34 — EL CAMPO MAGNÉTICO
% ==============================================================
% ==============================================================
\chapter{El Campo Magnético}

% ============================================================
\section{Fundamentos del Campo Magnético \texorpdfstring{$\vB{B}$}{B}}
% ============================================================

\subsection{Naturaleza e Historia del Magnetismo}

El magnetismo tiene su origen en la antigüedad clásica. El nombre proviene de
\textit{Magnesia}, región del Asia Menor donde se encontraron piedras naturales
llamadas \textbf{magnetita} ($\mathrm{Fe_3O_4}$) que atraían pequeños trozos de
hierro. Estas piedras son imanes naturales.

\textbf{Observaciones históricas fundamentales:}
\begin{itemize}[leftmargin=1.8em, itemsep=2pt]
  \item Los imanes naturales atraen hierro pero \textbf{no} oro ni plata:
        el magnetismo es selectivo por material.
  \item Dos imanes se atraen o repelen dependiendo de la orientación
        relativa de sus extremos.
  \item Una aguja magnetizada suspendida libremente se orienta siempre
        en la misma dirección geográfica (principio de la brújula).
\end{itemize}

\textbf{Aplicaciones modernas del magnetismo:}
\begin{itemize}[leftmargin=1.8em, itemsep=2pt]
  \item Grabación magnética (cintas, discos duros, tarjetas).
  \item Resonancia Magnética Nuclear (RMN) para imágenes médicas.
  \item Motores eléctricos y generadores.
  \item Medición del magnetismo planetario para estudiar estructuras internas.
  \item Aceleradores de partículas (ciclotrón, sincrotrón).
\end{itemize}

% ----------------------------------------------------------
\subsection{Definición y Representación del Campo Magnético \texorpdfstring{$\vB{B}$}{B}}
% ----------------------------------------------------------

Así como el espacio alrededor de una carga eléctrica está ocupado por un campo
eléctrico $\vB{E}$, el espacio alrededor de un imán o de un conductor con
corriente está ocupado por un \textbf{campo magnético}, representado por el
vector $\vB{B}$.

\begin{tcolorbox}[notebox, title=Nota terminológica importante]
  Existen dos vectores de campo en magnetismo:
  \begin{itemize}[leftmargin=1.5em, itemsep=1pt]
    \item $\vB{B}$: \textit{inducción magnética} o \textit{densidad de flujo
          magnético}. Es la cantidad \textbf{más fundamental}.
    \item $\vB{H}$: \textit{intensidad del campo magnético} (se usa en medios
          materiales).
  \end{itemize}
  En este texto $\vB{B}$ se denomina simplemente
  \textbf{``el campo magnético''}.
\end{tcolorbox}

% ----------------------------------------------------------
\subsection{Analogía entre Campo Eléctrico y Campo Magnético}
% ----------------------------------------------------------

En electrostática:
\[
  \text{carga eléctrica} \;\longleftrightarrow\; \vB{E}
  \;\longleftrightarrow\; \text{carga eléctrica}
\]

Una analogía tentadora sería:
\[
  \text{carga magnética} \;\longleftrightarrow\; \vB{B}
  \;\longleftrightarrow\; \text{carga magnética}
\]
Sin embargo, los \textit{monopolos magnéticos} no han sido hallados
experimentalmente. La relación operativa correcta es:

\begin{tcolorbox}[defbox, title=Relación fundamental del campo magnético]
  \[
    \text{corriente eléctrica} \;\longleftrightarrow\; \vB{B}
    \;\longleftrightarrow\; \text{corriente eléctrica}
  \]
  El campo magnético es \textbf{creado} por cargas en movimiento (corrientes)
  y \textbf{actúa} sobre cargas en movimiento (corrientes).
\end{tcolorbox}

% ----------------------------------------------------------
\subsection{Líneas de Campo Magnético}
% ----------------------------------------------------------

Las líneas de $\vB{B}$ siguen las mismas reglas de trazado que las de
$\vB{E}$: la tangente en cada punto da la \textbf{dirección} de $\vB{B}$,
y la densidad de líneas es proporcional a la \textbf{magnitud} de $\vB{B}$.

\begin{tcolorbox}[propbox, title=Diferencia fundamental respecto a las líneas de $\vB{E}$]
  \begin{itemize}[leftmargin=1.5em, itemsep=2pt]
    \item Las líneas de $\vB{E}$ son \textbf{abiertas}: nacen en cargas
          positivas y terminan en negativas.
    \item Las líneas de $\vB{B}$ forman \textbf{anillos cerrados}: nunca
          comienzan ni terminan. Esto refleja la ausencia de monopolos
          magnéticos.
    \item La fuerza eléctrica $\vB{F}=q\vB{E}$ es \textbf{paralela} a
          $\vB{E}$; la fuerza magnética
          $\vB{F}=q\vB{v}\times\vB{B}$ es \textbf{perpendicular} a
          $\vB{B}$.
  \end{itemize}
\end{tcolorbox}

\figplaceholder{34-1}{Limaduras de hierro sobre papel que cubre un imán de barra.
La distribución indica el patrón de líneas de campo magnético
(Fig.\,1 del texto, Cap.\,34).}

% ----------------------------------------------------------
\subsection{Polos Magnéticos y Campo Magnético Terrestre}
% ----------------------------------------------------------

En un imán de barra las líneas emergen del \textbf{polo Norte} y entran por el
\textbf{polo Sur}, formando siempre lazos cerrados a través del interior.

\begin{itemize}[leftmargin=1.8em, itemsep=2pt]
  \item Polos \textbf{opuestos} (N–S) $\Rightarrow$ se \textbf{atraen}.
  \item Polos \textbf{iguales} (N–N ó S–S) $\Rightarrow$ se \textbf{repelen}.
\end{itemize}

\begin{tcolorbox}[notebox, title=Campo magnético terrestre]
  El polo magnético de la Tierra en la región ártica es un polo magnético
  \textbf{Sur} (por eso atrae al polo Norte de la brújula). Cerca del Ecuador
  las líneas son casi paralelas a la superficie y apuntan del Sur geográfico
  hacia el Norte geográfico.
\end{tcolorbox}

\figplaceholder{34-2}{Líneas del campo magnético de un imán de barra, mostrando
los polos norte y sur y el carácter cerrado de las líneas
(Fig.\,5 del texto, Cap.\,34).}

% ============================================================
\section{La Fuerza Magnética sobre una Carga en Movimiento}
% ============================================================

\subsection{Procedimiento Experimental para Definir \texorpdfstring{$\vB{B}$}{B}}

Se proyecta una carga $q$ con velocidad $\vB{v}$ por distintos puntos y
direcciones; se registran los resultados:

\begin{enumerate}[label=\textbf{Exp.\,\arabic*.},
                  leftmargin=3.2em, itemsep=4pt]
  \item \textbf{Perpendicularidad.} $\vB{F}$ actúa siempre $\perp$ a
        $\vB{v}$, sea cual sea la dirección de $\vB{v}$.
  \item \textbf{Dependencia angular.} $F\propto\sin\varphi$, donde
        $\varphi$ es el ángulo entre $\vB{v}$ y la dirección de $\vB{B}$.
  \item \textbf{Proporcionalidad con $v$.} $F\propto v$.
  \item \textbf{Proporcionalidad con $q$.} $F\propto|q|$; se invierte si
        $q$ cambia de signo.
\end{enumerate}

\subsection{Definición Formal de \texorpdfstring{$\vB{B}$}{B}}

\begin{tcolorbox}[defbox, title=Definición operacional del campo magnético]
  \textbf{Dirección de $\vB{B}$:} dirección de $\vB{v}$ para la cual
  $F=0$.\\[5pt]
  \textbf{Magnitud de $\vB{B}$:}
  \begin{equation}
    \boxed{B = \frac{F_{\perp}}{qv}}
    \label{eq:defB}
  \end{equation}
  donde $F_\perp$ es la fuerza máxima cuando $\vB{v}\perp\vB{B}$.
\end{tcolorbox}

\subsection{Expresión Vectorial de la Fuerza Magnética}

\begin{tcolorbox}[defbox, title=Fuerza magnética — Forma escalar y vectorial]
  \textbf{Forma escalar:}
  \begin{equation}
    \boxed{F = qvB\sin\varphi}
    \label{eq:Fmagescalar}
  \end{equation}
  donde $0^\circ\le\varphi\le180^\circ$ es el ángulo más pequeño entre
  $\vB{v}$ y $\vB{B}$.\\[6pt]
  \textbf{Forma vectorial:}
  \begin{equation}
    \boxed{\vB{F} = q\,\vB{v}\times\vB{B}}
    \label{eq:Fmagvect}
  \end{equation}
\end{tcolorbox}

\textbf{Regla de la mano derecha} para $\vB{v}\times\vB{B}$: apuntar los
dedos en la dirección de $\vB{v}$, doblarlos hacia $\vB{B}$; el pulgar
señala $\vB{F}$ para $q>0$. Para $q<0$, $\vB{F}$ apunta en sentido
contrario.

\textbf{Casos extremos:}
\begin{itemize}[leftmargin=1.8em, itemsep=2pt]
  \item $\varphi=0^\circ$ ó $180^\circ$ $\Rightarrow$
        $F=0$ (sin deflexión).
  \item $\varphi=90^\circ$ $\Rightarrow$
        $F=F_{\max}=|q|vB$.
\end{itemize}

\figplaceholder{34-3}{Relación geométrica entre los vectores $\vB{F}$,
$\vB{v}$ y $\vB{B}$ para una carga positiva en movimiento
(Fig.\,4 del texto, Cap.\,34).}

\subsection{Propiedades Fundamentales de la Fuerza Magnética}

\begin{tcolorbox}[propbox, title=Propiedades esenciales]
  \begin{enumerate}[leftmargin=1.5em, itemsep=4pt]
    \item $\vB{F}\perp\vB{v}$ en todo instante.
    \item \textbf{La fuerza magnética no realiza trabajo:}
      \[
        dW = \vB{F}\cdot d\vB{s} = \vB{F}\cdot\vB{v}\,dt
           = 0 \qquad (\text{pues }\vB{F}\perp\vB{v})
      \]
    \item Por el teorema trabajo–energía: $\Delta K=0$, de modo que
      \[
        |\vB{v}| = \text{constante bajo }\ \vB{B}\ \text{ estático.}
      \]
    \item La fuerza es \textbf{deflectora}: curva la trayectoria sin
          cambiar la rapidez.
  \end{enumerate}
\end{tcolorbox}

\subsection{Unidades del Campo Magnético}

De la Ec.\,\eqref{eq:defB}:
\[
  [B] = \frac{\mathrm{N}}{\mathrm{C}\cdot\mathrm{m/s}}
       = \frac{\mathrm{N}}{\mathrm{A}\cdot\mathrm{m}}
\]

La unidad SI es el \textbf{tesla} (T):
\begin{equation}
  \boxed{1\un{T} = 1\,\frac{\mathrm{N}}{\mathrm{A}\cdot\mathrm{m}}
                 = 1\,\frac{\mathrm{kg}}{\mathrm{A}\cdot\mathrm{s}^2}}
\end{equation}

Relación con la unidad CGS: $1\un{T}=10^4\un{G}$ (gauss).

\vspace{6pt}
\begin{center}
\rowcolors{2}{LightNavy!50}{white}
\begin{tabular}{lc}
  \toprule
  \textbf{Ubicación} & $\boldsymbol{B}$ \textbf{(T)} \\
  \midrule
  Sala blindada magnéticamente & ${\sim}10^{-14}$ \\
  Espacio interestelar         & ${\sim}10^{-10}$ \\
  Superficie de la Tierra      & ${\sim}10^{-4}$  \\
  Imán pequeño de barra        & ${\sim}10^{-2}$  \\
  Electroimán grande de lab.   & ${\sim}1$        \\
  Imán superconductor          & ${\sim}5$        \\
  Estrella de neutrones (calc.)& ${\sim}10^{8}$   \\
  \bottomrule
\end{tabular}
\captionof{table}{Valores típicos del campo magnético (Tabla\,1, Cap.\,34).}
\end{center}

% ============================================================
\section{La Fuerza de Lorentz}
% ============================================================

\subsection{Definición}

Cuando actúan simultáneamente $\vB{E}$ y $\vB{B}$ sobre una partícula
cargada, la fuerza total es la \textbf{fuerza de Lorentz}:

\begin{equation}
  \boxed{\vB{F} = q\vB{E} + q\,\vB{v}\times\vB{B}}
  \label{eq:Lorentz}
\end{equation}

\begin{itemize}[leftmargin=1.8em, itemsep=2pt]
  \item Componente eléctrica $q\vB{E}$: actúa sobre \textit{cualquier}
        partícula cargada, en reposo o en movimiento.
  \item Componente magnética $q\vB{v}\times\vB{B}$: actúa
        \textit{únicamente} sobre partículas \textbf{en movimiento}.
\end{itemize}

\subsection{El Selector de Velocidad}

Con $\vB{E}\perp\vB{B}$ y ambos $\perp\vB{v}$, las fuerzas eléctrica y
magnética se oponen. La condición $\vB{F}_\text{neta}=\vB{0}$ exige:

\begin{equation}
  qE = qvB \implies
  \boxed{v = \frac{E}{B}}
  \label{eq:selectorv}
\end{equation}

Solo las partículas con exactamente $v=E/B$ pasan sin desviarse,
\textbf{independientemente de $q$ y $m$}.

\subsection{Experimento de J.\,J.\ Thomson — Medición de \texorpdfstring{$e/m$}{e/m} (1897)}

\begin{tcolorbox}[exbox, title=Procedimiento de Thomson]
  \textbf{Paso\,1 — Solo campo $E$:} la desviación vertical del haz es:
  \begin{equation}
    y = -\frac{qEL^2}{2mv^2}
    \label{eq:Thomson1}
  \end{equation}
  \textbf{Paso\,2 — Agregar $B$:} se ajusta hasta que la desviación es
  cero $\Rightarrow v=E/B$.\\[5pt]
  \textbf{Paso\,3 — Resultado:}
  \begin{equation}
    \boxed{\frac{e}{m} = \frac{2yE}{B^2L^2}}
    \label{eq:Thomson2}
  \end{equation}
  Thomson obtuvo $e/m\approx1.7\times10^{11}\un{C/kg}$.\\
  Valor moderno: $e/m=1.758\,819\,62\times10^{11}\un{C/kg}$.
\end{tcolorbox}

\figplaceholder{34-4}{Versión moderna del aparato de J.\,J.\,Thomson para medir
$e/m$ del electrón. Filamento F, placas P, bobinas de campo B, pantalla S
(Fig.\,8 del texto, Cap.\,34).}

\subsection{El Espectrómetro de Masas}

El radio de la trayectoria semicircular de un ion con velocidad $v=E/B$ en el
campo $B'$ es:
\[
  r = \frac{mv}{|q|B'} = \frac{mE}{|q|BB'}
\]
Despejando la razón carga–masa (espectrómetro Bainbridge):
\begin{equation}
  \boxed{\frac{q}{m} = \frac{E}{rBB'}}
  \label{eq:masspectrometer}
\end{equation}

\figplaceholder{34-5}{Diagrama del espectrómetro de masas Bainbridge. Selector de
velocidad + campo $B'$ que separa iones por su masa
(Fig.\,9 del texto, Cap.\,34).}

% ============================================================
\section{Cargas Circulantes en Campo Magnético Uniforme}
% ============================================================

\subsection{Condiciones para el Movimiento Circular}

Cuando $\varphi=90^\circ$ (velocidad perpendicular a $\vB{B}$) la fuerza
magnética tiene magnitud constante y es siempre $\perp\vB{v}$: condiciones
exactas para el \textbf{movimiento circular uniforme}.

\figplaceholder{34-6}{Electrones circulando en cámara al vacío con campo $\vB{B}$
perpendicular al plano. La trayectoria circular es visible por las
colisiones con el gas (Fig.\,10 del texto, Cap.\,34).}

\subsection{Radio de la Trayectoria Circular}

Igualando fuerza centrípeta y magnética:
\[
  \frac{mv^2}{r} = |q|vB
\]
\begin{equation}
  \boxed{r = \frac{mv}{|q|B} = \frac{p}{|q|B}}
  \label{eq:radiocircular}
\end{equation}
donde $p=mv$ es el \textit{ímpetu} lineal.

\begin{tcolorbox}[exbox,
  title=Ejemplo — Protón de $5.3\un{MeV}$ en $B=1.2\un{mT}$]
  \textbf{Velocidad} (a partir de $K=\tfrac{1}{2}mv^2$):
  \[
    v = \sqrt{\frac{2K}{m}}
      = \sqrt{\frac{2\times5.3\times1.60\times10^{-13}}{1.67\times10^{-27}}}
      = 3.2\times10^{7}\un{m/s}
  \]
  \textbf{Fuerza magnética} ($\varphi=90^\circ$):
  \[
    F = |q|vB
      = (1.60\times10^{-19})(3.2\times10^{7})(1.2\times10^{-3})
      = 6.1\times10^{-15}\un{N}
  \]
  \textbf{Dirección:} $\vB{v}$ al Norte, $\vB{B}$ hacia arriba $\Rightarrow$
  $\vB{F}$ apunta al \textbf{Este} (regla de la mano derecha).\\[4pt]
  \textit{Verificación relativista:}
  $K=5.3\un{MeV}\ll m_pc^2=938\un{MeV}$ $\checkmark$
\end{tcolorbox}

\subsection{Velocidad Angular y Frecuencia Ciclotrón}

\begin{equation}
  \boxed{\omega = \frac{|q|B}{m}}
  \label{eq:omega_ciclotron}
\end{equation}

\begin{equation}
  \boxed{\nu = \frac{\omega}{2\pi} = \frac{|q|B}{2\pi m}}
  \label{eq:nu_ciclotron}
\end{equation}

\begin{tcolorbox}[propbox, title=Isocronismo del movimiento ciclotrón]
  La frecuencia ciclotrón $\nu$ es \textbf{independiente de la velocidad}
  $v$. Una partícula lenta (círculo pequeño) y una rápida (círculo grande)
  tardan exactamente el mismo tiempo $T=1/\nu=2\pi m/(|q|B)$ en
  completar una vuelta.\\[4pt]
  \textbf{Condición de validez:} $v\ll c$ (régimen no relativista).
\end{tcolorbox}

% ============================================================
\section{Aceleradores de Partículas Cargadas}
% ============================================================

\subsection{El Ciclotrón}

\textbf{Descripción estructural:} dos objetos metálicos huecos en forma de
D (\textit{dees}) de cobre, con un entrehierro entre sus bordes rectos.
Un campo magnético $\vB{B}$ uniforme y perpendicular llena todo el espacio;
un oscilador eléctrico alterna la diferencia de potencial entre las dees.
En el centro, una fuente $S$ emite los iones a acelerar.

\figplaceholder{34-7}{Esquema del ciclotrón: dees D, fuente de iones S,
oscilador eléctrico y campo magnético $\vB{B}$ perpendicular al plano.
Trayectoria en espiral creciente (Fig.\,12 del texto, Cap.\,34).}

\textbf{Principio de operación:}
\begin{enumerate}[leftmargin=2.2em, itemsep=4pt]
  \item Los iones son \textbf{acelerados en el entrehierro} por el campo
        eléctrico oscilante.
  \item \textbf{Dentro de cada dee} (conductor hueco, $E=0$), el campo
        magnético curva la trayectoria en un semicírculo.
  \item Gracias al isocronismo, el ion llega al entrehierro exactamente
        cuando el campo se ha invertido (\textbf{condición de resonancia}).
  \item La partícula espirala hacia afuera ganando energía en cada paso.
\end{enumerate}

\textbf{Condición de resonancia:}
\begin{equation}
  \nu_\text{oscilador} = \nu_\text{ciclotrón} = \frac{|q|B}{2\pi m}
  \label{eq:resonancia}
\end{equation}

\textbf{Velocidad y energía al salir} (radio máximo $R$):
\begin{equation}
  \boxed{v_\text{final} = \frac{|q|BR}{m}}
  \label{eq:vfinal}
\end{equation}

\begin{equation}
  \boxed{K = \frac{q^2B^2R^2}{2m}}
  \label{eq:Kfinal}
\end{equation}

\begin{tcolorbox}[notebox, title=Observación sorprendente]
  La energía final $K$ depende de $B$ y $R$, pero \textbf{no} de la
  diferencia de potencial del oscilador. Una diferencia mayor acelera
  más por paso, pero la partícula da menos vueltas; la energía final
  es la misma.
\end{tcolorbox}

\begin{tcolorbox}[exbox,
  title=Ejemplo — Ciclotrón para deuterones ($R=75\un{cm}$,
  $B=1.5\un{T}$; $m_d=3.34\times10^{-27}\un{kg}$, $q=+e$)]
  \textit{(a) Frecuencia del oscilador:}
  \[
    \nu = \frac{(1.60\times10^{-19})(1.5)}
              {2\pi(3.34\times10^{-27})}
        = 1.1\times10^{7}\un{Hz} = 11\un{MHz}
  \]
  \textit{(b) Energía máxima:}
  \[
    K = \frac{(1.60\times10^{-19})^2(1.5)^2(0.75)^2}
            {2(3.34\times10^{-27})}
      = 4.85\times10^{-12}\un{J} = 30\un{MeV}
  \]
\end{tcolorbox}

\subsection{Corrección Relativista — Límite del Ciclotrón}

Para $v$ comparable a $c$, el ímpetu relativista es
$p=mv/\!\sqrt{1-v^2/c^2}$, y la frecuencia ciclotrón se vuelve
dependiente de $v$:

\begin{equation}
  \nu = \frac{|q|B\sqrt{1-v^2/c^2}}{2\pi m}
  \label{eq:nu_rel}
\end{equation}

La condición de resonancia \textbf{se pierde} para energías altas
($\gtrsim50\un{MeV}$ para protones, donde $v/c\approx0.3$).

\subsection{El Sincrotrón}

En el sincrotrón, $B$ y $\nu$ se \textbf{varían} sincrónicamente para
mantener el radio de la órbita constante. Se usan muchos imanes individuales
distribuidos a lo largo de un gran anillo de radio fijo.

\begin{tcolorbox}[exbox, title=Sincrotrón del Fermilab (referencia histórica)]
  \begin{itemize}[leftmargin=1.5em, itemsep=2pt]
    \item Circunferencia: $\approx6.4\un{km}$ (4 millas).
    \item Número de imanes: $\approx3\,000$.
    \item Los protones dan $\sim400\,000$ vueltas para alcanzar
          $\mathbf{1\,000\un{GeV}=1\un{TeV}}$.
    \item Duración del proceso: $\sim10\un{s}$ por ráfaga.
  \end{itemize}
\end{tcolorbox}

% ============================================================
\section{Confinamiento Magnético y Fenómenos Asociados}
% ============================================================

\subsection{El Espejo Magnético}

Una partícula con componente de velocidad paralela a $\vB{B}$ describe una
trayectoria \textbf{helicoidal}. En una \textit{botella magnética} (campo
más intenso en los extremos), la componente longitudinal se desacelera y la
partícula es \textbf{reflejada}, quedando atrapada entre los dos extremos.

\figplaceholder{34-8}{Trayectoria helicoidal de una partícula cargada en una
botella magnética. Las componentes centrípetas en los extremos reflejan
la partícula (Fig.\,14 del texto, Cap.\,34).}

\textbf{Aplicación:} confinamiento de plasmas calientes en investigación de
fusión termonuclear (tokamaks).

\subsection{Los Cinturones de Radiación Van Allen}

El campo magnético terrestre actúa como espejo magnético a escala planetaria:

\begin{itemize}[leftmargin=1.8em, itemsep=2pt]
  \item \textbf{Cinturón interno:} principalmente protones atrapados.
  \item \textbf{Cinturón externo:} principalmente electrones atrapados.
  \item Las partículas oscilan en espiral entre los polos en tiempos de
        segundos y derivan lentamente alrededor de la Tierra.
\end{itemize}

Descubiertos por James Van Allen en 1958 con los primeros satélites
artificiales estadounidenses.

\subsection{Cálculo Numérico de Trayectorias — Método de Euler}

Para $\vB{B}=B_0\hat{z}$, las ecuaciones de movimiento son:
\[
  F_x = qv_yB_0,\qquad F_y = -qv_xB_0,\qquad F_z=0
\]

Integración numérica con paso temporal $\delta t$:
\begin{align}
  \delta v_x &= \frac{qB_0}{m}\,v_y\,\delta t \\[4pt]
  \delta v_y &= -\frac{qB_0}{m}\,v_x\,\delta t
\end{align}

Actualización con velocidad promedio:
\[
  x_\text{nuevo} = x + \tfrac{1}{2}(v_{x,\text{viejo}}+v_{x,\text{nuevo}})\,\delta t
\]

Para campo no uniforme
$B=B_0\!\left(1+\tfrac{x^2+y^2}{R^2}\right)$, la trayectoria forma un
patrón en forma de \textbf{flor}.

% ============================================================
\section{El Efecto Hall}
% ============================================================

\subsection{Principio Físico}

\begin{tcolorbox}[notebox, title=Contexto histórico]
  Descubierto por \textbf{Edwin H.\ Hall} en 1879, estudiante de posgrado
  de 24 años en la Universidad Johns Hopkins — \textbf{veinte años antes}
  del descubrimiento del electrón por Thomson (1897).
\end{tcolorbox}

Una cinta de anchura $w$ y espesor $t$ conduce corriente $i$; un campo
$\vB{B}$ perpendicular desvía los portadores de carga lateralmente,
acumulándolos en un lado y generando un campo eléctrico transversal
$E_H$ (voltaje Hall $V_H$).

\figplaceholder{34-9}{Cinta conductora en campo $\vB{B}$: (a)\,situación
al activarse $B$; (b)\,equilibrio. Los portadores negativos
se acumulan en el lado derecho (Fig.\,18 del texto, Cap.\,34).}

\textbf{Condición de equilibrio} (fuerza de Lorentz neta = 0):
\begin{align}
  q\vB{E}_H + q\vB{v}_d\times\vB{B} &= \vB{0} \notag\\
  E_H &= v_d B \label{eq:HallEq}
\end{align}

\subsection{Densidad de Portadores de Carga}

Con la relación $v_d=i/(newt)$:

\begin{equation}
  \boxed{V_H = \frac{iB}{net}}
  \label{eq:VHall}
\end{equation}

\begin{equation}
  \boxed{n = \frac{iB}{etV_H}}
  \label{eq:nHall}
\end{equation}

\textbf{Resistencia Hall:}
\begin{equation}
  R_H = \frac{V_H}{i} = \frac{B}{net}
  \label{eq:RHall}
\end{equation}

\subsection{Signo de los Portadores y Datos Experimentales}

El signo de $V_H$ determina si los portadores son electrones (negativo) o
huecos (positivo):

\vspace{6pt}
\begin{center}
\rowcolors{2}{LightNavy!50}{white}
\begin{tabular}{lccc}
  \toprule
  \textbf{Material} &
  $\boldsymbol{n}\ (10^{28}/\mathrm{m}^3)$ &
  \textbf{Signo de }$V_H$ &
  \textbf{Portadores/átomo} \\
  \midrule
  Na              & 2.5                & $-$ & 0.99              \\
  K               & 1.5                & $-$ & 1.1               \\
  Cu              & 11                 & $-$ & 1.3               \\
  Ag              & 7.4                & $-$ & 1.3               \\
  Al              & 21                 & $-$ & 3.5               \\
  Sb              & 0.31               & $-$ & 0.09              \\
  Be              & 2.6                & $+$ & 2.2               \\
  Zn              & 19                 & $+$ & 2.9               \\
  Si (puro)       & $1.5\times10^{-12}$& $-$ & $3\times10^{-13}$ \\
  Si (tipo\,\textit{n}) & $10^{-7}$   & $-$ & $2\times10^{-8}$  \\
  \bottomrule
\end{tabular}
\captionof{table}{Datos del efecto Hall (Tabla\,2, Cap.\,34).}
\end{center}

\begin{tcolorbox}[exbox,
  title=Ejemplo — Cinta de Cu ($t=150\,\mu\mathrm{m}$,
  $B=0.65\un{T}$, $i=23\un{A}$; $n=8.49\times10^{28}\un{m}^{-3}$)]
  \[
    V_H = \frac{iB}{net}
        = \frac{(23)(0.65)}
               {(8.49\times10^{28})(1.60\times10^{-19})(150\times10^{-6})}
        = 7.3\times10^{-6}\un{V} = 7.3\,\mu\mathrm{V}
  \]
\end{tcolorbox}

\subsection{El Efecto Hall Cuantizado (von Klitzing, 1980)}

A temperaturas $\sim1\un{K}$ y campos de varios Tesla, en muestras
bidimensionales (heterostructuras GaAs/AlGaAs), $R_H$ exhibe escalones
planos con valores exactamente:

\begin{equation}
  \boxed{R_H = \frac{h}{je^2},\quad j=1,2,3,\ldots}
  \label{eq:HallQ}
\end{equation}

Primer escalón ($j=1$): $R_H=h/e^2=25\,812.806\un{\Omega}\approx25.8\un{k\Omega}$.

\begin{tcolorbox}[propbox, title=Importancia metrológica — Premio Nobel 1985]
  $h/e^2=25\,812.806\un{\Omega}$ se conoce con precisión mejor que
  $1$ parte en $10^{10}$. Desde 1990 es el \textbf{estándar internacional
  del ohm}.\\[4pt]
  \textbf{Klaus von Klitzing} obtuvo el Nobel de Física 1985 por este
  descubrimiento.\\[4pt]
  En 1982 se descubrió el \textbf{Efecto Hall Cuantizado Fraccionario}
  (Tsui, Störmer; Nobel 1998), con escalones en fracciones de $e^2$.
\end{tcolorbox}

% ============================================================
\section{La Fuerza Magnética sobre una Corriente}
% ============================================================

\subsection{Origen Físico y Deducción}

Los electrones de conducción son desviados lateralmente por
$\vB{F}=q\vB{v}\times\vB{B}$; al estar confinados en el conductor,
esta fuerza se transmite al conductor sólido.

Para $N=nAL$ electrones en un segmento de longitud $L$,
cargando $q=-e$, con velocidad de arrastre $\vB{v}_d$:
\[
  \vB{F} = -nALe\,\vB{v}_d\times\vB{B}
\]

Definiendo el vector $\vB{L}$ (magnitud $L$, dirección de la corriente)
y usando $-nALe\vB{v}_d=i\vB{L}$:

\begin{equation}
  \boxed{\vB{F} = i\,\vB{L}\times\vB{B}}
  \label{eq:Fcorriente}
\end{equation}

Para elemento diferencial:
\begin{equation}
  d\vB{F} = i\,d\vB{s}\times\vB{B}
  \label{eq:dFcorriente}
\end{equation}

\figplaceholder{34-10}{Alambre flexible entre polos de imán: (a)\,sin corriente,
(b)\,con corriente, (c)\,corriente invertida. La desviación cambia de
sentido (Fig.\,20 y 21 del texto, Cap.\,34).}

\begin{tcolorbox}[exbox,
  title=Ejemplo — Alambre en equilibrio ($m/L=46.6\un{g/m}$,
  $i=28\un{A}$)]
  \[
    iLB = mg \implies B = \frac{(m/L)g}{i}
            = \frac{(46.6\times10^{-3})(9.8)}{28}
            = 1.6\times10^{-2}\un{T} = 16\un{mT}
  \]
  Este campo es $\approx400$ veces el campo magnético terrestre.
\end{tcolorbox}

\begin{tcolorbox}[propbox, title=Resultado notable — Trayectorias en campo uniforme]
  La fuerza sobre \textit{cualquier} trayectoria conductora entre dos puntos
  en un campo \textbf{uniforme} es igual a la fuerza sobre el segmento
  \textbf{recto} que une esos puntos.
\end{tcolorbox}

% ============================================================
\section{Momento de Torsión sobre una Espira de Corriente}
% ============================================================

\subsection{Configuración y Cálculo del Torque}

Espira rectangular de lados $a$ y $b$ suspendida en campo uniforme
$\vB{B}$. El vector normal $\hat{n}$ forma ángulo $\theta$ con $\vB{B}$.

\figplaceholder{34-11}{Espira rectangular en campo uniforme $\vB{B}$. Fuerzas
sobre los cuatro lados y momento de torsión resultante. Eje de rotación
en la dirección $z$ (Fig.\,26 del texto, Cap.\,34).}

La fuerza neta sobre la espira es \textbf{cero} ($\sum\vB{F}=\vB{0}$);
no hay traslación. El torque es:

\begin{equation}
  \boxed{\tau = NiAB\sin\theta}
  \label{eq:torque_espira}
\end{equation}

donde $A=ab$ y $N$ es el número de vueltas.
\textbf{Dirección:} el torque tiende a alinear $\hat{n}$ con $\vB{B}$
($\theta\to0$). La Ec.\,\eqref{eq:torque_espira} es válida para
\textit{cualquier} espira plana.

\subsection{Aplicación: El Galvanómetro}

Campo magnético radial $\Rightarrow$ $\theta=90^\circ$ siempre.
Equilibrio entre torque magnético y resorte restaurador:
\[
  NiAB = \kappa\phi \implies i = \frac{\kappa}{NAB}\,\phi
\]

La deflexión $\phi$ es proporcional a $i$ (escala calibrada).

\figplaceholder{34-12}{Esquema del galvanómetro: bobina giratoria, imán permanente
con núcleo de hierro dulce para campo radial, resorte restaurador y aguja
(Fig.\,27 del texto, Cap.\,34).}

\begin{tcolorbox}[exbox,
  title=Ejemplo — Galvanómetro (250 vueltas, $B=0.23\un{T}$,
  $i=100\,\mu\mathrm{A}$, $\phi=28^\circ=0.49\un{rad}$)]
  \[
    A = (0.021)(0.012) = 2.52\times10^{-4}\un{m}^2
  \]
  \[
    \kappa = \frac{NiAB}{\phi}
           = \frac{(250)(100\times10^{-6})(2.52\times10^{-4})(0.23)}{0.49}
           = 3.0\times10^{-6}\un{N\cdot m/rad}
  \]
\end{tcolorbox}

% ============================================================
\section{El Dipolo Magnético}
% ============================================================

\subsection{Definición del Momento Dipolar Magnético}

\begin{equation}
  \boxed{\mu = NiA} \qquad [\mathrm{A\cdot m^2} = \mathrm{J/T}]
  \label{eq:mu}
\end{equation}

Dirección: paralela a $\hat{n}$ (regla de la mano derecha aplicada a la
corriente en la espira).

\subsection{Torque y Energía Potencial del Dipolo}

\begin{equation}
  \boxed{\vB{\tau} = \vB{\mu}\times\vB{B}}
  \label{eq:tau_dipolo}
\end{equation}

\begin{equation}
  \boxed{U = -\vB{\mu}\cdot\vB{B} = -\mu B\cos\theta}
  \label{eq:U_dipolo}
\end{equation}

Análogas a $\vB{\tau}=\vB{p}\times\vB{E}$ y $U=-\vB{p}\cdot\vB{E}$
del dipolo eléctrico.

\textbf{Equilibrios:}
\begin{itemize}[leftmargin=1.8em, itemsep=2pt]
  \item $\theta=0^\circ$: $U=-\mu B$ (mínimo) $\Rightarrow$
        equilibrio \textbf{estable}.
  \item $\theta=180^\circ$: $U=+\mu B$ (máximo) $\Rightarrow$
        equilibrio \textbf{inestable}.
\end{itemize}

\textbf{Trabajo para girar $180^\circ$:}
\[
  W = \Delta U = 2\mu B
\]

\subsection{Valores de Momentos Dipolares Magnéticos}

\begin{center}
\rowcolors{2}{LightNavy!50}{white}
\begin{tabular}{lc}
  \toprule
  \textbf{Sistema} & $\boldsymbol{\mu}$ \textbf{(J/T)} \\
  \midrule
  Núcleo del átomo de H        & $2.0\times10^{-28}$ \\
  El protón                    & $1.4\times10^{-26}$ \\
  El electrón                  & $9.3\times10^{-24}$ \\
  El átomo de nitrógeno        & $2.8\times10^{-23}$ \\
  Bobina pequeña típica        & $5.4\times10^{-6}$  \\
  Imán de barra pequeño        & $5$                 \\
  Bobina superconductora       & $400$               \\
  La Tierra                    & $8.0\times10^{22}$  \\
  \bottomrule
\end{tabular}
\captionof{table}{Momentos dipolares magnéticos (Tabla\,3, Cap.\,34).}
\end{center}

\begin{tcolorbox}[propbox,
  title=Consecuencias de la comparación nuclear vs.\ atómica]
  Los momentos nucleares son $3$–$6$ órdenes de magnitud \textbf{menores}
  que los atómicos:
  \begin{enumerate}[leftmargin=1.8em, itemsep=2pt]
    \item Los electrones \textbf{no pueden} ser constituyentes del núcleo.
    \item Los efectos magnéticos ordinarios de los materiales están
          gobernados por el magnetismo \textbf{atómico} (electrónico).
    \item Para alinear dipolos nucleares (RMN) se requieren campos
          $3$–$6$ órdenes de magnitud \textbf{mayores} que para dipolos
          atómicos a igual temperatura.
  \end{enumerate}
\end{tcolorbox}

% ============================================================
\section*{Resumen de Ecuaciones — Capítulo 34}
\addcontentsline{toc}{section}{Resumen de ecuaciones — Capítulo\,34}
% ============================================================

\begin{center}
\renewcommand{\arraystretch}{1.7}
\rowcolors{2}{LightNavy!50}{white}
\begin{longtable}{p{0.38\textwidth}p{0.52\textwidth}}
  \toprule
  \textbf{Ecuación} & \textbf{Descripción} \\
  \midrule
  $\displaystyle\vB{F}=q\,\vB{v}\times\vB{B}$ &
        Fuerza magnética sobre carga en movimiento \\
  $\displaystyle\vB{F}=q\vB{E}+q\,\vB{v}\times\vB{B}$ &
        Fuerza de Lorentz \\
  $\displaystyle B=F_\perp/(qv)$ &
        Definición operacional de $B$ \\
  $\displaystyle r=mv/(|q|B)=p/(|q|B)$ &
        Radio de trayectoria circular \\
  $\displaystyle\omega=|q|B/m$ &
        Velocidad angular ciclotrón \\
  $\displaystyle\nu=|q|B/(2\pi m)$ &
        Frecuencia ciclotrón (no relativista) \\
  $\displaystyle K=q^2B^2R^2/(2m)$ &
        Energía máxima en ciclotrón \\
  $\displaystyle v=E/B$ &
        Velocidad en selector de velocidad \\
  $\displaystyle n=iB/(etV_H)$ &
        Densidad de portadores (efecto Hall) \\
  $\displaystyle R_H=h/(je^2)$ &
        Efecto Hall cuantizado \\
  $\displaystyle\vB{F}=i\,\vB{L}\times\vB{B}$ &
        Fuerza sobre segmento de corriente \\
  $\displaystyle d\vB{F}=i\,d\vB{s}\times\vB{B}$ &
        Fuerza sobre elemento diferencial \\
  $\displaystyle\tau=NiAB\sin\theta$ &
        Torque sobre espira (escalar) \\
  $\displaystyle\vB{\tau}=\vB{\mu}\times\vB{B}$ &
        Torque sobre espira (vectorial) \\
  $\displaystyle\mu=NiA$ &
        Momento dipolar magnético \\
  $\displaystyle U=-\vB{\mu}\cdot\vB{B}$ &
        Energía potencial del dipolo magnético \\
  \bottomrule
\end{longtable}
\end{center}

\clearpage

% ==============================================================
% ==============================================================
%  CAPÍTULO 35 — LA LEY DE AMPÈRE
% ==============================================================
% ==============================================================
\chapter{La Ley de Amp\`{e}re}

% ============================================================
\section{El Experimento de Oersted y la Fuente del Campo Magnético}
% ============================================================

\subsection{Descubrimiento de Oersted (1820)}

Hans Christian Oersted observó que una aguja de brújula colocada cerca de un
alambre recto por el que pasaba una corriente se orientaba
\textbf{perpendicular} al alambre, independientemente de la orientación del
mismo. Este fue el primer vínculo experimental entre electricidad y magnetismo.

\figplaceholder{35-1}{El experimento de Oersted: brújulas colocadas alrededor de
un alambre recto con corriente. La aguja se orienta perpendicular al alambre
(Fig.\,1 del texto, Cap.\,35).}

\textbf{Interpretación moderna:} la corriente en el alambre crea un campo
magnético $\vB{B}$ que ejerce un torque sobre la aguja de la brújula y la
alinea con ese campo.

\subsection{Analogía Metodológica con la Electrostática}

\begin{center}
\renewcommand{\arraystretch}{1.5}
\begin{tabular}{p{0.42\textwidth}p{0.42\textwidth}}
  \toprule
  \textbf{Electrostática} & \textbf{Magnetostática} \\
  \midrule
  Ley de Coulomb: campo de elemento de carga $dq$:
  $d\vB{E}=\dfrac{1}{4\pi\varepsilon_0}\dfrac{dq}{r^2}\hat{u}_r$
  &
  Ley de Biot-Savart: campo de elemento de corriente $ids$:
  $d\vB{B}=\dfrac{\mu_0}{4\pi}\dfrac{ids\times\hat{u}_r}{r^2}$
  \\[8pt]
  Ley de Gauss (simetría alta): $\oint\vB{E}\cdot d\vB{A}=q_\text{enc}/\varepsilon_0$
  &
  Ley de Ampère (simetría alta): $\oint\vB{B}\cdot d\vB{s}=\mu_0 i$
  \\
  \bottomrule
\end{tabular}
\end{center}

% ============================================================
\section{La Ley de Biot-Savart}
% ============================================================

\subsection{La Constante de Permeabilidad Magnética del Vacío}

La constante de proporcionalidad en la ley de la fuerza magnética se
define fijando:

\begin{equation}
  k = \frac{\mu_0}{4\pi} = 10^{-7}\un{T\cdot m/A}
\end{equation}

donde $\mu_0$ es la \textbf{constante de permeabilidad magnética del vacío}:

\begin{equation}
  \boxed{\mu_0 = 4\pi\times10^{-7}\un{T\cdot m/A}
        \approx 1.2566\times10^{-6}\un{T\cdot m/A}}
  \label{eq:mu0}
\end{equation}

Esta constante define la unidad de corriente (el ampere) a través de la ley
de la fuerza entre conductores paralelos. Se relaciona con $\varepsilon_0$
a través de la velocidad de la luz:
\[
  c = \frac{1}{\sqrt{\mu_0\varepsilon_0}}
\]

\subsection{Enunciado Formal de la Ley de Biot-Savart}

\begin{tcolorbox}[defbox, title=Ley de Biot-Savart]
  El campo magnético $d\vB{B}$ en un punto $P$ creado por un elemento de
  corriente $i\,d\vB{s}$ a una distancia $r$ es:
  \begin{equation}
    \boxed{d\vB{B} = \frac{\mu_0}{4\pi}\,
      \frac{i\,d\vB{s}\times\hat{u}_r}{r^2}
      = \frac{\mu_0}{4\pi}\,
      \frac{i\,d\vB{s}\times\vB{r}}{r^3}}
    \label{eq:BiotSavart}
  \end{equation}
  donde $\vB{r}$ apunta del elemento de corriente al punto $P$ y
  $\hat{u}_r=\vB{r}/r$.\\[5pt]
  \textbf{Magnitud:}
  \begin{equation}
    dB = \frac{\mu_0}{4\pi}\,\frac{i\,ds\sin\theta}{r^2}
    \label{eq:BSmag}
  \end{equation}
  donde $\theta$ es el ángulo entre $d\vB{s}$ y $\vB{r}$.\\[5pt]
  \textbf{Campo total:}
  \begin{equation}
    \vB{B} = \frac{\mu_0}{4\pi}\int\frac{i\,d\vB{s}\times\hat{u}_r}{r^2}
    \label{eq:BSintegral}
  \end{equation}
\end{tcolorbox}

\figplaceholder{35-2}{Geometría de la ley de Biot-Savart: elemento de corriente
$i\,d\vB{s}$, vector $\vB{r}$ al punto $P$, y dirección de $d\vB{B}$
perpendicular al plano de $d\vB{s}$ y $\vB{r}$ (Fig.\,4 del texto, Cap.\,35).}

% ============================================================
\section{Aplicaciones de la Ley de Biot-Savart}
% ============================================================

\subsection{Campo de un Alambre Recto Largo}

\figplaceholder{35-3}{Alambre recto largo con corriente $i$: geometría para
la integración de la ley de Biot-Savart. El elemento $i\,dx$ y las
variables $r$, $\theta$, $R$ (Fig.\,5 del texto, Cap.\,35).}

Tomando $x$ como variable de integración a lo largo del alambre,
con $r=\sqrt{x^2+R^2}$ y $\sin\theta=R/\sqrt{x^2+R^2}$:

\[
  B = \frac{\mu_0 i}{4\pi}\int_{-\infty}^{+\infty}
      \frac{R\,dx}{(x^2+R^2)^{3/2}}
    = \frac{\mu_0 i}{4\pi}\cdot\frac{2}{R}
\]

\begin{equation}
  \boxed{B = \frac{\mu_0 i}{2\pi R}}
  \label{eq:Balamlargo}
\end{equation}

El campo disminuye como $1/R$ con la distancia $R$ al alambre.

\subsection{Campo sobre el Eje de un Anillo Circular de Corriente}

Para un anillo de radio $R$ con corriente $i$, a distancia axial $z$
del centro:

Por simetría, las componentes transversales de $d\vB{B}$ se cancelan;
solo sobrevive la componente axial:

\begin{equation}
  \boxed{B = \frac{\mu_0 i R^2}{2(R^2+z^2)^{3/2}}}
  \label{eq:Banillo}
\end{equation}

\figplaceholder{35-4}{Anillo circular de corriente $i$ de radio $R$. Cálculo
del campo axial en el punto $P$ a distancia $z$. Componentes
$d\vB{B}_\parallel$ y $d\vB{B}_\perp$ (Fig.\,6 del texto, Cap.\,35).}

\textbf{Casos especiales:}
\begin{itemize}[leftmargin=1.8em, itemsep=2pt]
  \item Centro del anillo ($z=0$):
    $\displaystyle B = \frac{\mu_0 i}{2R}$
  \item Punto lejano ($z\gg R$):
    $\displaystyle B \approx \frac{\mu_0 i R^2}{2z^3}
                           = \frac{\mu_0}{2\pi}\,\frac{\mu}{z^3}$
    (campo dipolar con $\mu=iA$)
\end{itemize}

\subsection{Tabla Comparativa: Dipolos Eléctrico y Magnético}

\begin{center}
\renewcommand{\arraystretch}{1.6}
\rowcolors{2}{LightNavy!50}{white}
\begin{tabular}{p{0.30\textwidth}p{0.28\textwidth}p{0.28\textwidth}}
  \toprule
  \textbf{Propiedad} &
  \textbf{Dipolo eléctrico} &
  \textbf{Dipolo magnético} \\
  \midrule
  Torque en campo externo &
  $\vB{\tau}=\vB{p}\times\vB{E}$ &
  $\vB{\tau}=\vB{\mu}\times\vB{B}$ \\
  Energía en campo externo &
  $U=-\vB{p}\cdot\vB{E}$ &
  $U=-\vB{\mu}\cdot\vB{B}$ \\
  Campo axial ($z\gg$) &
  $\displaystyle E=\frac{p}{2\pi\varepsilon_0 z^3}$ &
  $\displaystyle B=\frac{\mu_0\mu}{2\pi z^3}$ \\
  Campo ecuatorial ($x\gg$) &
  $\displaystyle E=\frac{p}{4\pi\varepsilon_0 x^3}$ &
  $\displaystyle B=\frac{\mu_0\mu}{4\pi x^3}$ \\
  \bottomrule
\end{tabular}
\captionof{table}{Analogía entre dipolo eléctrico y dipolo magnético
(Tabla\,1, Cap.\,35).}
\end{center}

\subsection{Campo de una Cinta Plana Conductora}

Para una cinta de anchura $a$ y corriente $i$, el campo en el punto $P$
sobre la bisectriz perpendicular a distancia $R$:

\begin{equation}
  B = \frac{\mu_0 i}{\pi a}\arctan\!\left(\frac{a}{2R}\right)
  \label{eq:Bcinta}
\end{equation}

Para $R\gg a$: $B\approx\mu_0 i/(2\pi R)$ (recupera el resultado del
alambre delgado).

% ============================================================
\section{Líneas del Campo Magnético}
% ============================================================

Las líneas de $\vB{B}$ cerca de un alambre recto forman
\textbf{círculos concéntricos} con espaciamiento que crece como $1/r$.
La dirección se determina por la regla de la mano derecha: el pulgar en
la dirección de la corriente; los demás dedos encorvan en la dirección
de $\vB{B}$.

\figplaceholder{35-5}{Líneas del campo magnético (círculos concéntricos)
alrededor de un alambre recto largo con corriente
(Fig.\,9 del texto, Cap.\,35).}

\begin{tcolorbox}[notebox, title=Principio para calcular fuerzas]
  La fuerza magnética sobre un conductor se calcula usando únicamente el
  \textbf{campo externo} en que está inmerso. El campo que el propio
  conductor genera no ejerce fuerza sobre sí mismo (igual que la gravedad
  terrestre no actúa sobre la Tierra misma).
\end{tcolorbox}

% ============================================================
\section{Interacción entre Conductores Paralelos}
% ============================================================

\subsection{Fuerza entre Dos Alambres Paralelos}

\figplaceholder{35-6}{Dos alambres paralelos con corrientes $i_1$ e $i_2$
en la misma dirección. El campo $\vB{B}_1$ en el alambre 2 y la fuerza
$\vB{F}_{21}$ de atracción (Fig.\,11 del texto, Cap.\,35).}

El alambre 1 crea en la posición del alambre 2 (separación $d$):
\[
  B_1 = \frac{\mu_0 i_1}{2\pi d}
\]

La fuerza sobre longitud $L$ del alambre 2:
\begin{equation}
  \boxed{F_{21} = \frac{\mu_0 L\, i_1 i_2}{2\pi d}}
  \label{eq:Fparalelos}
\end{equation}

\begin{tcolorbox}[propbox, title=Regla general de interacción entre corrientes]
  \begin{itemize}[leftmargin=1.5em, itemsep=2pt]
    \item Corrientes \textbf{paralelas} (mismo sentido)
          $\Rightarrow$ se \textbf{atraen}.
    \item Corrientes \textbf{antiparalelas} (sentidos opuestos)
          $\Rightarrow$ se \textbf{repelen}.
  \end{itemize}
  Esta regla es \textit{opuesta} a la de las cargas eléctricas
  (cargas iguales se repelen, pero corrientes iguales se atraen).
\end{tcolorbox}

\subsection{Definición del Ampere}

\begin{tcolorbox}[defbox, title=Definición SI del ampere]
  Dados dos alambres paralelos, de sección transversal circular
  despreciable, separados $1\un{m}$ en el vacío: se define
  $\mathbf{1\un{ampere}}$ como la corriente en cada alambre que
  produce una fuerza de $\mathbf{2\times10^{-7}\un{N}}$ por metro
  de longitud.
\end{tcolorbox}

% ============================================================
\section{La Ley de Amp\`{e}re}
% ============================================================

\subsection{Motivación — Analogía con la Ley de Gauss}

La ley de Biot-Savart permite calcular $\vB{B}$ para cualquier
distribución de corriente, pero la integración puede ser compleja.
De manera análoga a la ley de Gauss para campos eléctricos, existe
una ley que aprovecha la \textbf{simetría} para simplificar el
cálculo de $\vB{B}$:

\begin{tcolorbox}[defbox, title=Ley de Amp\`{e}re]
  \begin{equation}
    \boxed{\oint\vB{B}\cdot d\vB{s} = \mu_0\, i_\text{enc}}
    \label{eq:Ampere}
  \end{equation}
  donde $i_\text{enc}$ es la corriente neta \textbf{encerrada} por
  la trayectoria cerrada (anillo amperiano), y la integral de línea
  se evalúa a lo largo de dicho anillo.
\end{tcolorbox}

\textbf{La ley de Ampère es una de las cuatro ecuaciones de Maxwell.}
Es más fundamental que la ley de Biot-Savart.

\subsection{El Anillo Amperiano — Procedimiento}

Análogo a la superficie gaussiana, se construye una
\textbf{curva cerrada imaginaria} (anillo amperiano) bien elegida:

\begin{enumerate}[leftmargin=2.2em, itemsep=4pt]
  \item Elegir el anillo de modo que $\vB{B}$ sea
        \textbf{constante en magnitud} sobre la trayectoria,
        y \textbf{tangente} a ella (o perpendicular).
  \item La integral se simplifica: $\oint B\,ds=B\cdot(\text{longitud})$.
  \item La corriente encerrada $i_\text{enc}$ solo incluye las corrientes
        que \textbf{atraviesan} la superficie delimitada por el anillo.
\end{enumerate}

\textbf{Convención de signos:} regla de la mano derecha — si los dedos se
doblan en el sentido de recorrido del anillo, el pulgar indica la
dirección positiva de la corriente.

\subsection{Corrientes Externas al Anillo}

Las corrientes externas contribuyen al campo $\vB{B}$ en los puntos
sobre el anillo, pero sus contribuciones a $\oint\vB{B}\cdot d\vB{s}$
se cancelan exactamente al integrar alrededor de la trayectoria cerrada.

\figplaceholder{35-7}{Ley de Ampère aplicada a anillo arbitrario que encierra
corrientes $i_1$ e $i_2$ y excluye $i_3$. Se muestran el anillo amperiano
y las corrientes (Fig.\,13 del texto, Cap.\,35).}

\subsection{Aplicación — Campo del Alambre Recto Largo}

Anillo amperiano: círculo de radio $r$ coaxial con el alambre.

Por simetría: $\vB{B}$ constante y tangente en todo el anillo.
$\theta=0^\circ$ en toda la trayectoria.

\[
  \oint B\,ds = B(2\pi r) = \mu_0 i
\]

\begin{equation}
  \boxed{B = \frac{\mu_0 i}{2\pi r}}
  \label{eq:Bext_alambre}
\end{equation}

Idéntico al resultado de Biot-Savart, obtenido con mucho menos esfuerzo.

\figplaceholder{35-8}{Anillo amperiano circular de radio $r$ para el alambre recto
largo. $\vB{B}$ es tangente y constante en todo el anillo
(Fig.\,14 del texto, Cap.\,35).}

\subsection{Aplicación — Campo Dentro de un Conductor Cilíndrico Sólido}

Para $r<R$ (dentro del alambre de radio $R$ con corriente $i$ uniforme):

\[
  B(2\pi r) = \mu_0 i\,\frac{\pi r^2}{\pi R^2}
\]

\begin{equation}
  \boxed{B = \frac{\mu_0 i\,r}{2\pi R^2}} \qquad (r\le R)
  \label{eq:Bint_alambre}
\end{equation}

\begin{tcolorbox}[propbox, title=Campo magnético dentro y fuera del alambre]
  \begin{itemize}[leftmargin=1.5em, itemsep=2pt]
    \item Para $r<R$ (interior): $B\propto r$ (crece linealmente).
    \item Para $r=R$ (superficie): ambas expresiones coinciden.
    \item Para $r>R$ (exterior): $B\propto 1/r$ (decrece).
    \item El campo \textbf{máximo} ocurre en la superficie $r=R$.
  \end{itemize}
\end{tcolorbox}

\figplaceholder{35-9}{Gráfica de $B$ en función de $r$ para un alambre cilíndrico
sólido de radio $R$. Crecimiento lineal interior, decaimiento como $1/r$
exterior (Fig.\,16 del texto, Cap.\,35).}

% ============================================================
\section{Solenoides y Toroides}
% ============================================================

\subsection{El Solenoide}

El solenoide es un alambre largo devanado en hélice compacta, que conduce
corriente $i_0$. Su longitud es mucho mayor que su diámetro.

\figplaceholder{35-10}{Sección de solenoide ``extendido'' mostrando líneas de
campo magnético. Campo uniforme interior, campo casi nulo exterior
(Fig.\,17 y 18 del texto, Cap.\,35).}

\textbf{Campo externo:} para un solenoide ideal infinito, $B_\text{ext}=0$
(los campos de las espiras superiores e inferiores se cancelan en el exterior).

\textbf{Aplicación de la Ley de Ampère — Anillo amperiano rectangular $abcd$:}

Solo el segmento interior $ab$ contribuye ($Bh$). Los demás segmentos dan
contribución nula por perpendicularidad o por $B=0$ exterior.

Con $n$ espiras por unidad de longitud y corriente $i_0n h$ encerrada:

\begin{equation}
  Bh = \mu_0 i_0 n h \implies
  \boxed{B = \mu_0 i_0 n}
  \label{eq:Bsolenoide}
\end{equation}

\begin{tcolorbox}[propbox, title=Propiedades del campo en el solenoide ideal]
  \begin{itemize}[leftmargin=1.5em, itemsep=2pt]
    \item $B$ es \textbf{uniforme} en toda la sección transversal interior.
    \item $B$ es independiente del diámetro y la longitud del solenoide.
    \item $B$ depende únicamente de $i_0$ y $n$ (espiras/m).
    \item $B_\text{exterior}\approx0$ (solenoide práctico largo).
  \end{itemize}
\end{tcolorbox}

\figplaceholder{35-11}{Anillo amperiano rectangular $abcd$ en el solenoide ideal.
Los segmentos $bc$ y $da$ son perpendiculares a $\vB{B}$; el segmento $cd$
está en la región exterior donde $B=0$ (Fig.\,19 del texto, Cap.\,35).}

\begin{tcolorbox}[exbox,
  title=Ejemplo — Solenoide ($L=1.23\un{m}$, $d=3.55\un{cm}$,
  5 capas de 850 espiras, $i_0=5.57\un{A}$)]
  \[
    n = \frac{5\times850}{1.23\un{m}} = 3\,455\un{espiras/m}
  \]
  \[
    B = \mu_0 i_0 n
      = (4\pi\times10^{-7})(5.57)(3\,455)
      = 2.42\times10^{-2}\un{T} = 24.2\un{mT}
  \]
\end{tcolorbox}

\subsection{El Toroide}

El toroide es un solenoide doblado en forma de rosca de radio interior
$a$ y radio exterior $b$.

\figplaceholder{35-12}{Toroide con $N$ espiras. Anillo amperiano circular de radio
$r$ interior para calcular el campo magnético
(Fig.\,20 del texto, Cap.\,35).}

Anillo amperiano: círculo de radio $r$ concéntrico con el eje del toroide.
La corriente total encerrada es $Ni_0$.

\begin{equation}
  B(2\pi r) = \mu_0 i_0 N \implies
  \boxed{B = \frac{\mu_0 i_0 N}{2\pi r}} \qquad (a\le r\le b)
  \label{eq:Btoroide}
\end{equation}

\begin{itemize}[leftmargin=1.8em, itemsep=2pt]
  \item $B$ \textbf{no} es uniforme en la sección transversal del toroide:
        varía con $r$.
  \item $B=0$ en el exterior ($r>b$ ó $r<a$) del toroide ideal.
  \item Para $a\approx b$ (toroide muy delgado), el denominador $2\pi r$
        es la circunferencia media y $N/(2\pi r)=n$ (espiras/m), de modo
        que $B\approx\mu_0 i_0 n$: mismo que el solenoide.
\end{itemize}

\textbf{Aplicación:} el toroide es el componente central del
\textit{tokamak}, reactor de fusión termonuclear.

% ============================================================
\section{Electromagnetismo y Marcos de Referencia}
% ============================================================

\subsection{Análisis en el Marco de Reposo $S$}

Una partícula de carga $q>0$ en reposo junto a un alambre con corriente $i$
(Fig.\,22a del texto):

\begin{itemize}[leftmargin=1.8em, itemsep=2pt]
  \item El alambre es eléctricamente neutro: $\lambda_+=|\lambda_-|$,
        campo eléctrico neto $\vB{E}=\vB{0}$.
  \item Existe campo $\vB{B}\ne\vB{0}$, pero la partícula está en reposo
        $\Rightarrow$ no hay fuerza magnética.
  \item \textbf{Conclusión en $S$:} fuerza neta $= 0$.
\end{itemize}

\subsection{Análisis en el Marco $S'$ (en Movimiento a $v_d$)}

El marco $S'$ se mueve paralelo al alambre a la velocidad de arrastre $v_d$
de los electrones:

\begin{itemize}[leftmargin=1.8em, itemsep=2pt]
  \item En $S'$: los electrones están en reposo; los iones se mueven.
  \item La partícula cargada se mueve $\Rightarrow$ experimenta fuerza
        magnética $F_B$ alejándose del alambre.
  \item Por la \textbf{contracción de Lorentz}: la barra de iones
        (en movimiento en $S'$) queda más corta $\Rightarrow$
        $\lambda'_+>|\lambda'_-|$.
  \item Hay densidad de carga neta positiva en $S'$ $\Rightarrow$
        campo eléctrico $F_E$ que repele a $q$.
  \item Cálculo detallado: $F_E$ cancela exactamente a $F_B$
        $\Rightarrow$ fuerza neta $=0$ en $S'$ también.
\end{itemize}

\begin{tcolorbox}[propbox,
  title=Conclusión fundamental — Relatividad del electromagnetismo]
  Los campos eléctrico y magnético \textbf{no son entidades independientes}.
  Un campo puramente eléctrico en un marco tiene componentes magnéticas en
  otro, y viceversa. Las leyes de Maxwell son
  \textbf{invariantes bajo la transformación de Lorentz}.\\[5pt]
  \textit{Einstein (1905):} ``La fuerza sobre un cuerpo en movimiento
  dentro de un campo magnético no es otra cosa que un campo eléctrico.''\\[5pt]
  El \textbf{magnetismo} puede verse como un
  \textbf{efecto relativista} de la electricidad,
  observable a velocidades mucho menores que la de la luz debido a la
  enorme concentración de cargas en los conductores.
\end{tcolorbox}

% ============================================================
\section*{Resumen de Ecuaciones — Capítulo 35}
\addcontentsline{toc}{section}{Resumen de ecuaciones — Capítulo\,35}
% ============================================================

\begin{center}
\renewcommand{\arraystretch}{1.7}
\rowcolors{2}{LightNavy!50}{white}
\begin{longtable}{p{0.40\textwidth}p{0.50\textwidth}}
  \toprule
  \textbf{Ecuación} & \textbf{Descripción} \\
  \midrule
  $\displaystyle\mu_0=4\pi\times10^{-7}\un{T\cdot m/A}$ &
        Constante de permeabilidad del vacío \\
  $\displaystyle c=1/\!\sqrt{\mu_0\varepsilon_0}$ &
        Velocidad de la luz en función de $\mu_0$ y $\varepsilon_0$ \\
  $\displaystyle d\vB{B}=\frac{\mu_0}{4\pi}
        \frac{i\,d\vB{s}\times\hat{u}_r}{r^2}$ &
        Ley de Biot-Savart \\
  $\displaystyle B=\frac{\mu_0 i}{2\pi R}$ &
        Campo del alambre recto largo (a distancia $R$) \\
  $\displaystyle B=\frac{\mu_0 iR^2}{2(R^2+z^2)^{3/2}}$ &
        Campo axial del anillo circular \\
  $\displaystyle B=\frac{\mu_0 i}{2R}$ &
        Campo en el centro del anillo ($z=0$) \\
  $\displaystyle F=\frac{\mu_0 L\,i_1i_2}{2\pi d}$ &
        Fuerza entre conductores paralelos \\
  $\displaystyle\oint\vB{B}\cdot d\vB{s}=\mu_0 i_\text{enc}$ &
        Ley de Ampère \\
  $\displaystyle B=\frac{\mu_0 i}{2\pi r}$ ($r>R$) &
        Campo exterior del alambre cilíndrico \\
  $\displaystyle B=\frac{\mu_0 i\,r}{2\pi R^2}$ ($r<R$) &
        Campo interior del alambre cilíndrico \\
  $\displaystyle B=\mu_0 i_0 n$ &
        Campo del solenoide ideal ($n$\,=\,espiras/m) \\
  $\displaystyle B=\frac{\mu_0 i_0 N}{2\pi r}$ &
        Campo del toroide ideal \\
  \bottomrule
\end{longtable}
\end{center}

% ============================================================
%  TABLA FINAL COMPARATIVA
% ============================================================
\section*{Tabla Comparativa General — Electricidad y Magnetismo}
\addcontentsline{toc}{section}{Tabla comparativa — Electricidad y Magnetismo}

\begin{center}
\renewcommand{\arraystretch}{1.6}
\rowcolors{2}{LightGuinda!50}{white}
\begin{tabular}{p{0.30\textwidth}p{0.30\textwidth}p{0.30\textwidth}}
  \toprule
  \textbf{Concepto} &
  \textbf{Electricidad} &
  \textbf{Magnetismo} \\
  \midrule
  Fuente del campo &
  Carga $q$ &
  Corriente $i$ ó $q$ en mov. \\
  Vector campo &
  $\vB{E}$ &
  $\vB{B}$ \\
  Constante fundamental &
  $\varepsilon_0=8.85\times10^{-12}\un{C^2/N\cdot m^2}$ &
  $\mu_0=4\pi\times10^{-7}\un{T\cdot m/A}$ \\
  Ley de la fuerza &
  Coulomb: $F=kq_1q_2/r^2$ &
  Biot-Savart: $d\vB{B}=\frac{\mu_0}{4\pi}\frac{i\,d\vB{s}\times\hat{r}}{r^2}$\\
  Ley de simetría &
  Gauss: $\oint\vB{E}\cdot d\vB{A}=q/\varepsilon_0$ &
  Ampère: $\oint\vB{B}\cdot d\vB{s}=\mu_0 i$ \\
  Líneas de campo &
  Abiertas (cargas $\to$ cargas) &
  Cerradas (anillos) \\
  Dipolo &
  $\vB{p}$: par $\pm q$ &
  $\vB{\mu}=Ni A\hat{n}$: espira \\
  Torque &
  $\vB{\tau}=\vB{p}\times\vB{E}$ &
  $\vB{\tau}=\vB{\mu}\times\vB{B}$ \\
  Energía potencial &
  $U=-\vB{p}\cdot\vB{E}$ &
  $U=-\vB{\mu}\cdot\vB{B}$ \\
  \bottomrule
\end{tabular}
\captionof{table}{Comparación sistemática entre campo eléctrico y magnético.}
\end{center}

% ============================================================
\end{document}
% ============================================================%% It is just an empty TeX file.
%% Write your code here.